\section{Setup: a decomposition that answers when it should not}
\label{sec:intro}

A tissue model reports a differentiation marker down 90\% between a patient's
healthy and diseased sample. Two mechanisms produce that observation and they
imply different interventions. Either the cells making the marker have left, a
\emph{compositional} change, or those cells remain and produce less, an
\emph{intrinsic} change.
\iffull
Bulk measurement cannot separate them. Single-cell
measurement can, but only where enough of the relevant cells survive in both
arms.
\else
Separating them needs single-cell measurement, and only works where enough of
the relevant cells survive in both arms.
\fi
The separating estimator is Kitagawa's demographic standardisation
\citep{kitagawa1955}. Write $f$ for the fraction of cells of the type in
question and $m$ for their mean per-cell output:
\begin{equation}
\underbrace{(f_T - f_N)\, m_N}_{\text{compositional}} \;+\;
\underbrace{f_N\, (m_T - m_N)}_{\text{intrinsic}} \;+\;
\underbrace{(f_T - f_N)(m_T - m_N)}_{\text{interaction}}
\;=\; f_T m_T - f_N m_N .
\label{eq:kitagawa}
\end{equation}
Economists know this as Oaxaca--Blinder \citep{oaxaca1973,blinder1973},
epidemiologists as standardisation \citep{kline2011}. The same identity sits
under synthetic control arms \citep{abadie2010} and target trial emulation
\citep{hernan2016}. Any system reporting a group-resolved quantity inherits its
identifiability problem.

The problem is that Equation~\ref{eq:kitagawa} returns a number for any input.
The intrinsic term is a difference of two means over cells of one type. When the
diseased arm holds none of those cells, ``mean output per surviving cell of this
type'' has no referent, and the estimator still returns a number. So our
estimator returns one of three verdicts: an estimate, a wide-interval estimate,
or \notest{}, where the term is \texttt{None} and never $0.0$. Zero and ``we did
not ask'' are values a downstream average cannot tell apart. The rule gating
those verdicts is a cutpoint on $n$, the count of relevant cells in the diseased
arm: answer at $n \geq 50$, flag a wide interval from 20 to 49, abstain below 20.

\iffull
\paragraph{Related work.} Simulation-study methodology (ADEMP and its
forerunners), the differential-abundance methods that never face this
estimand, and the selective-prediction literature the abstention rule belongs
to are surveyed in Appendix~\ref{app:relatedwork}.

\paragraph{Relation to known circularity.} That a simulation study can favour
methods resembling its generator is not new, and we are not claiming it. Meng
\citep{meng1994} named the condition \emph{congeniality}; the plasmode
literature warns that synthetic benchmarks can be ``biased towards favoring
causal inference methods that mimic the modeling choices made when generating
the synthetic datasets'' \citep{franklin2014}; Curth et al.~\citep{curth2021}
show semi-synthetic benchmarks systematically favouring some estimators inside
this community, and Gentzel et al.~\citep{gentzel2019} argue for interventional
measures over simulated ones for the same reason; Shaw et al.~\citep{shaw2026} document the mirror image, a
generator that makes good estimators look bad; and in inverse problems the same
practice has had a name for two decades, the \emph{inverse crime}
\citep{wirgin2004,kaipio2007}. Simulation-based calibration
\citep{talts2018} is the closest existing instrument, and its own recent
critique \citep{modrak2025} is this paper's argument in another setting:
implementations in wide use ``could never detect large classes of problems''.

\textbf{What we add is that the degenerate case is a kind, not a degree, and
that this reverses the practical advice.} All of the above describe a bias of
unknown magnitude, and a bias predicts a curve that looks \emph{suspiciously
good} --- something a careful reader might catch. The case here is an identity:
the residual is exactly zero, the estimator is algebraically absent from the
ratio, and the curve therefore looks \emph{exactly as a correct one should},
near 1 and tightening with $n$. \textbf{A validation that fails loudly is a
nuisance; one that succeeds silently is a hazard}, and inspection cannot
separate them. So the contribution is not the phenomenon. It is the column ---
compare against the realised draw, and require the difference not to be
identically zero --- together with the finding that two estimators from
different literatures both fail it, and that
\S\ref{sec:calibration} then contradicts what the blind curve had certified.

\else
\paragraph{Related work.} Filters of this shape are standard rather than novel,
and are surveyed with the simulation-study methodology in
Appendix~\ref{app:relatedwork}. Our claim is not the filter but that it is a
selective-prediction rule \citep{chow1970,geifman2017} --- something you can
\emph{calibrate} and get wrong.
\fi

